\section{Crisis Staff Introductions}

\begin{chairletter}[Chair]{Peter Zakhar}
Hey all! My name is Peter, a (not so recent) graduate from Hult International Business School, originally from Slovakia. I completed my BBA program back in May 2025, while having run HultMUN as President and was the USG for ER at LIMUN 2025. During my time in MUN at university, I have attended almost 30 conferences from across the UK circuit, spanning Spain, France, Philippines, Taiwan, Peru, Thailand and Indonesia. My most recent triumph was leading HultMUN as a Faculty Advisor to win Best Large Delegation at Harvard WorldMUN’26. I’m also the Executive Director of Harvard WorldMUN’27 in Latvia (you are all welcome to join!) while full-time running a new MUN Platform, Gavelling, which I’m sure that I won’t shut up about 🙂

This will be my fourth time chairing IEUMUN, having started in 2023 and quickly falling in love with the city, the people and the conference in itself. My expectations for this IEUMUN, just like in other conferences I have attended and chairing, is exactly pointed to that personal development. I hope all delegates in the committees will get comfortable to explore their weak sides while emphasising their strong ones. It is only through that I managed to raise over 200 delegates personally with my training programs in MUN, while conducting a very welcoming and fun environment. Trust me, you’ll see that while the committee might be tough, we will be having a blast together! Prepare your research, because curveballs will be hitting the committee from every side.

\end{chairletter}

\begin{chairletter}[Co-Chair]{Miriam Liberal}
Hi! My name is Miriam and I'm born and raised in Spain. Despite most of my family being from Caceres (a rural area in Spain) I've lived all my life in Madrid and consider myself a true “gata” (Madrid’s  demonym). For all of those travelling for this MUN to Madrid, I sincerely hope you have time to explore this beautiful city, should you need any city recommendations don't hesitate to ask me.

For me MUN represents an opportunity to learn in practical environments and truly get a taste of the globalisation in which our world currently operates. I wish for all of us to imagine even further in this crisis, to get nervous, even confused and ultimately thrilled about what we are about to create during the committee. It is exactly that “thrilling” that got me into crisis committees, the feeling of not knowing what is about to happen mimicking what is real life. Apart from MUN I am a big fan of museums and interactive expositions (again, if you need any recommendations I would love to talk about my favourite places in the city), some of my favourite topics of research are culture, fashion, animals and  environment. As a last fun fact about me, I am a big fan of music and musicals, I have been singing for the most part of my life. I can’t wait to meet all of you, and hear your proposal in the committee. See you soon!

\end{chairletter}

\begin{chairletter}[Senior Backroomer]{Faris Zebib}
Hiya, I’m Faris, one of your head backroomers for this year. I’ve been Crisis Director for conferences across Europe, and organized many MUN’s through my own organization, BrusselsMUN, over the last couple of years.

IEUMUN has a special place in my heart, as it was one of the first conferences I ever did during its first edition post-Covid. This year will be a special one, with 4 times the committees, and 4 times the chaos!

Together, we will lead your cabinet to victory, but this game is a two-way street. So play your cards right, do your research, and let your creative juices flow!

I’ll catch you on the battlefield.

\end{chairletter}

\begin{chairletter}[Junior Backroomer]{Nina Maniurka}
Hi everyone! My name is Nina and I'm your Backroomer for this year's IEUMUN. I'm Polish, though I now split my time between Madrid, where I study Economics at UC3M, and Barcelona, where I'm currently interning. I've also studied at LSE, Bocconi and the University of California along the way, so I'm always happy to swap stories about student life across a few different countries.

I gravitated to crisis because backrooming is basically a puzzle that fights back. Every directive you send in changes the board, and my job is to make sure the board changes right back. The committees I've enjoyed most are the ones where delegates stop playing it safe and actually test how far they can push things, so consider that an open invitation.

Whether this is your first crisis or your tenth, I hope you'll use these three days to take a few risks you wouldn't normally take. I can't wait to meet you all and see what chaos we create together. See you in Madrid!

\end{chairletter}

\begin{chairletter}[Junior Backroomer]{Luka Hornstein Tomic}
Hi everyone! My name is Luka and I am your Junior Backroomer for this MUN. I am born and raised in Croatia with German and Bosnian roots. I speak both German and Croatian and am keen on learning as many languages as I can. I study and live in Madrid at IE and hope you will truly experience the city and the atmosphere in its true form during these 3 days.

This is my first time backrooming and I am so excited. I truly believe MUN is the best thing you could do and get interested in as a student as it gives you so much new knowledge, skills and new friendships. Exactly this is what you can expect from this years IEUMUN and I promise you it will be one to remember. I wish you all the best and to really step out your comfort zone because that’s what crisis is all about. Besides that a fun fact about me is that I love music in all forms from electronic to the classical sounds of an old instrument. I love playing sports and opening myself up to new experiences. I cant wait to meet you all and read your directives. All the best in committee!

\end{chairletter}
